% !TeX root = ../main.tex
Pooled CRISPR screens increasingly combine time courses, ordered FACS
fractions, donor cohorts, dose series, and factorial treatment arms.  These
experiments ask about slopes, adjusted associations, and interactions.
Reducing them to repeated high-versus-low comparisons discards intermediate
observations and prevents estimation of one effect conditional on the other
design variables.

The count denominator determines how guide abundance is modeled.
Negative-binomial methods generally model marginal counts and account for
library size through normalization or an offset
\citep{marioni2008rnaseq,anders2010deseq,love2014deseq2,lun2016edger}.
In a pooled screen, each guide count is observed together with the total number
of mapped guide reads in its library.  Conditioning on that total treats guide
abundance as a proportion.  A beta-binomial model represents this sampling
structure and separates depth-dependent sequencing variation from
heterogeneity among biological libraries.  The original
CB\textsuperscript{2} study reported better goodness of fit for beta-binomial
than negative-binomial models in pooled CRISPR counts \citep{jeong2019beta}.

CB\textsuperscript{2} applies this principle to two independent groups.
Beta-binomial regression is available in \texttt{corncob} and statistical
packages including \texttt{aod}, \texttt{VGAM}, and \texttt{gamlss}
\citep{martin2020corncob,lesnoff2012aod,yee1996vgam,rigby2005gamlss}, but these
tools lack a CRISPR-oriented workflow connecting full-library totals,
guide-level regression, and screen-level outputs.  Established screen and
RNA-sequencing methods, including MAGeCK-MLE \citep{li2014mageck}, edgeR-QL
\citep{lun2016edger}, DESeq2 \citep{love2014deseq2}, and limma--voom
\citep{law2014voom}, already support multivariable designs.  Chronos and
Waterbear provide specialized models for population dynamics and joint
FACS-bin membership, respectively
\citep{dempster2021chronos,pimentel2024waterbear}.  These methods provide
comparators for assessing the practical value of extending the beta-binomial
screen model.

We developed BARCS (Beta-binomial Analysis and Regression for CRISPR Screens)
to carry the library-total-conditional model into multivariable analysis.
BARCS replaces two group means with a regression design matrix, allowing a
single fit to estimate time, donor-adjusted abundance, treatment, or interaction
effects.  Independent biological libraries determine the residual degrees of
freedom.  Guide-level estimation is followed by optional dispersion moderation,
gene aggregation, and control-based calibration.  The default denominator
retains full-library totals through guide filtering; a separate control-based
normalization option estimates effects relative to a chosen control class.

We first describe the BARCS workflow and evaluate its longitudinal and
covariate-adjusted coefficients in Cas13 and IL2RA screens.  Genome-scale and
factorial simulations then test dispersion moderation and denominator choice
against known truth.  A null-calibration grid examines the consequences of
combining correlated guides, and an external audit assesses the effect of
pre-fit subsetting on a published false-discovery result.  These analyses
establish the uses and limitations of multivariable beta-binomial regression,
including the distinction between recovering biological signals and obtaining
calibrated guide- or gene-level probabilities.
